would not express the strain in the body correctly, and the terms of
(27) that contain \(p, q, r\) would not represent a displacement
possible in a rigid body. Bearing this restriction in mind, we conclude
that, if the six components of strain are given, the corresponding
displacement is arbitrary to the extent of an additional displacement of
the type expressed by (27); but, if we impose 
six independent conditions,
such as that, at the origin, the displacement \((u, v, w)\) and the rotation
\(\left(\varpi_x, \varpi_y, \varpi_z\right)\) vanish, or again that, at the
same point
\begin{equation}\label{eq:28}
u=0,v=0,w=0, \frac{\partial u}{\partial z}=0,
\frac{\partial v}{\partial z}=0, \frac{\partial v}{\partial x}=0,
\end{equation}
the expression for the displacement with given strains will be unique.
The particular set of equations (\ref{eq:28}) indicate that one point of the body
(the origin), one linear element of the body (that along the axis of
\(z\) issuing from the origin) and one plane-element of the body (that
in the plane of \(z, x\) containing the origin) retain their positions
after the strain. It is manifestly possible, after straining a body in
any way, to bring it back by translation and rotation so that a given
point, a given linear element through the point and a given
plane-element through the line shall recover their primitive positions.

\section{Curvilinear orthogonal coordinates}
For many problems it is convenient to use systems of curvilinear
coordinates\footnote{The theory is due to Lamé. See his
Leçons sur les coordonnées curvilignes, Paris, 1859.}
instead of the ordinary Cartesian coordinates. These may be
introduced as follows: Let \(f(x, y, z)=\alpha \), some constant,
be the
equation of a surface. If \(\alpha \) is allowed to vary we obtain a
family of surfaces. In general one surface of the family will pass
through a chosen point, and a neighbouring point will in general lie
on
a neighbouring surface of the family, so that \(\alpha \) is a function
of \(x, y, z\), viz., the function denoted by \(f\). If
\(\alpha+d \alpha \) is the parameter of that surface of the family which
passes through \((x+d x, y+d y, z+d z)\), we have
\[
d \alpha=\frac{\partial f}{\partial x} d x+\frac{\partial f}{\partial y} d y+
\frac{\partial f}{\partial z} d z=\frac{\partial \alpha}{\partial x} d x+
\frac{\partial \alpha}{\partial y} d y+\frac{\partial \alpha}{\partial z} d z.
\]

If we have three independent families of surfaces given by the equations
\[f_1(x, y, z)=\alpha, \quad f_2(x, y, z)=\beta, \quad f_3(x, y, z)=\gamma,\]
so that in general one surface of each family passes through a chosen
point, then a point may be determined by the values of
\(\alpha, \beta, \gamma \) which belong to the surfaces that pass through
it\footnote{The determination of the
point may not be free from ambiguity, e.g., in elliptic coordinates, an
ellipsoid and two confocal hyperboloids pass through any point, and they
meet in seven other points. The ambiguity is removed if the region of
space considered is suitably limited, e g., in the case of elliptic
coordinates, if it is an octant bounded by principal planes.},
and a neighbouring point will be determined by the neighbouring values
\(\alpha+d \alpha, \beta+d \beta, \gamma+d \gamma \). Such quantities as
\(\alpha, \beta, \gamma \) are called ``curvilinear coordinates'' of the
point.

The most convenient systems of curvilinear coordinates for applications
to the theory of Elasticity are determined by families of surfaces which
cut each other everywhere at right angles. In such a case we have a
triply-orthogonal family of surfaces. It is well known that there exists
an infinite number of sets of such surfaces, and, according to a
celebrated theorem due to Dupin, the line of intersection of two
surfaces belonging to different families of such a set is a line of
curvature on each\footnote{Salmon, Geometry of three
dimensions, 4th ed., p.269}. In what follows we shall take
\(\alpha, \beta, \gamma \) to be the parameters of such a set of
surfaces, so that the following relations hold:
\[
\langle \nabla \beta, \nabla \gamma \rangle =
\langle \nabla \gamma, \nabla \alpha \rangle =
\langle \nabla \alpha, \nabla \beta \rangle = 0.
\]

The length of the normal, \(d n_1\), to a surface of the family \(\alpha \)
intercepted between the surfaces \(\alpha \) and \(\alpha+d \alpha \) is
determined by the observation that the direction-cosines of the normal
to \(\alpha \) at the point \((x, y, z)\) are
\begin{equation}\label{eq:29}
\frac{1}{h_1} \frac{\partial \alpha}{\partial x}, \quad \frac{1}{h_1}
\frac{\partial \alpha}{\partial y}, \quad \frac{1}{h_1} \frac{\partial
\alpha}{\partial z},
\end{equation}
where \(h_1\) is expressed by the first of equations (31) below. For, by
projecting the line joining two neighbouring points on the normal to
\(\alpha \), we obtain the equation
\begin{equation}\label{eq:30}
d n_1= \frac{1}{h_1}\left(\frac{\partial \alpha}{\partial x} d x+
\frac{\partial \alpha}{\partial y} d y+
\frac{\partial \alpha}{\partial z} d z\right)=\frac{d \alpha}{h_1}
\end{equation}

In like manner the elements \(d n_2, d n_3\) of the normals to \(\beta \)
and \(\gamma \) are \(d \beta / h_2\) and \(d \gamma / h_3\), where
\begin{equation}\label{eq:31}
\left.\begin{array}{l}
h_1^2={\left(\frac{\partial \alpha}{\partial x}\right)}^2+
{\left(\frac{\partial \alpha}{\partial y}\right)}^2+
{\left(\frac{\partial \alpha}{\partial z}\right)}^2,\\
h_2^2=
{\left(\frac{\partial \beta}{\partial x}\right)}^2+
{\left(\frac{\partial \beta}{\partial y}\right)}^2+
{\left(\frac{\partial \beta}{\partial z}\right)}^2,\\
h_3^2=
{\left(\frac{\partial \gamma}{\partial x}\right)}^2+
{\left(\frac{\partial \gamma}{\partial y}\right)}^2+
{\left(\frac{\partial \gamma}{\partial z}\right)}^2.
\end{array}\right\}
\end{equation}
\todo[inline]{This is page 52. The
printing of the third edition of this equation
incorrectly replaces \(h_3^2\) with a second \(h_2^2\).
The correct equation appears in the second and fourth editions,
and also in every other printing of the third edition
that I have checked.}

The distance between two neighbouring points being
\({\left(d n_1{ }^2+d n_2{ }^2+d n_3{ }^2\right)}^{\frac{1}{2}}\), we have
the expression for the ``line-element'' \(d s\), i.e.\ the distance
between the points \((\alpha, \beta, \gamma)\) and
\((\alpha+d \alpha, \beta+d \beta, \gamma+d \gamma)\), in the form
\[
{(d s)}^2={\left(d \alpha / h_1\right)}^2+
{\left(d \beta / h_2\right)}^2+{\left(d \gamma / h_3\right)}^2.
\]
In general \(h_1, h_2, h_3\) are regarded as functions of
\(\alpha, \beta, \gamma \).

\section[Curvilinear Strain]{Components of strain referred to curvilinear orthogonal coordinates}
%\section{Components of strain referred to curvilinear orthogonal coordinates} %20

\todo[inline]{Love put The next footnote in the section title.}

Anchor\footnote{The method here given is due to Borchardt, J.  f. Math.
(Crelle), Bd. 76 (1873).  Another method is given in the `Note on applications
of moving axes' at the end of this book.}

The length in the unstrained state of the line joining the points \((\alpha,
\beta, \gamma)\) and \((\alpha+d \alpha, \beta+d \beta, \gamma+d \gamma)\) is
given by (32); we seek the length in the strained state of the line joining the
same pair of particles. Let \(u_a, u_\beta, u_\gamma \) be the projections of
the displacement of any particle on the normals to the surfaces \(\alpha,
\beta, \gamma \) that pass through its position in the unstrained state. When
the displacement is small the coordinates of the point occupied by a particle
are changed from \(\alpha, \beta, \gamma \) to
\(\alpha+h_1 u_\alpha, \beta+h_2 u_\beta, \gamma+h_3 u_\gamma \). The
\(\alpha \)-coordinate \((\alpha+d \alpha)\) of a neighbouring point is changed into
\[
\alpha+d \alpha+\left\{h_1 u_\alpha+d \alpha \frac{\partial}{\partial
\alpha}\left(h_1 u_\alpha \right)+d \beta \frac{\partial}{\partial
\beta}\left(h_1 u_\alpha \right)+d \gamma \frac{\partial}{\partial
\gamma}\left(h_1 u_\alpha \right)\right\}.
\]
and similar changes are made in the \(\beta \) and \(\gamma \)-coordinates
(\(\beta+d \beta \) and \(\gamma+d \gamma \)).  Again, the values of
\( \{h_i\} \) at a displaced particle differ from those at its undisplaced
position. For example, \(1 / h_1\) is changed into
\[
\frac{1}{h_1}+h_2 u_\alpha \frac{\partial}{\partial
\alpha}\left(\frac{1}{h_1}\right)+h_2 u_\beta \frac{\partial}{\partial
\beta}\left(\frac{1}{h_1}\right)+h_3 u_\gamma \frac{\partial}{\partial
\gamma}\left(\frac{1}{h_1}\right).
\]
It follows that \(d \alpha / h_1\) must be replaced by
\[
\begin{aligned} {\left[ d \alpha
\left(1+\frac{\partial}{\partial \alpha} \left(h_1 u_\alpha \right)\right)+
d \beta \frac{\partial}{\partial \beta}\left(h_2 u_\alpha \right)+
d \gamma \frac{\partial}{\partial \gamma}\left(h_1 u_\alpha \right)\right]}
\times \\
\times\left[ \frac{1}{h_1}+
h_1 u_\alpha \frac{\partial}{\partial \alpha} \left(\frac{1}{h_1}\right)+
h_2 u_\beta \frac{\partial}{\partial \beta}\left(\frac{1}{h_1}\right)+
h_3 u_\gamma \frac{\partial}{\partial \gamma}\left(\frac{1}{h_1}\right)
\right],
\end{aligned}
\]
and, when products of quantities of the order \(u_a\) are neglected, this is
\[
\begin{aligned}
\frac{d \alpha}{h_1}+d \alpha \left\{\frac{\partial u_\alpha}{\partial
\alpha}+h_2 u_\beta \frac{\partial}{\partial
\beta}\left(\frac{1}{h_1}\right)\right. & \left.+h_3 u_\gamma
\frac{\partial}{\partial \gamma}\left(\frac{1}{h_1}\right)\right\} \\
& +d \beta \frac{1}{h_1} \frac{\partial}{\partial \beta}\left(h_1
u_\alpha \right)+d \gamma \frac{1}{h_1} \frac{\partial}{\partial
\gamma}\left(h_1 u_\alpha \right) .
\end{aligned}
\]
Similar changes must be made in \(d \beta / h_2\) and
\(d \gamma / h_3\). The length of the line joining two particles in the
strained state is found by forming the square root of the sum of the
squares of the three expressions of the type (33). Let \(d s\) be the
length of the linear element in the unstrained state, and let
\(l, m, n\) be the direction-cosines of it referred to the normals to
the surfaces \(\alpha, \beta, \gamma \) at a point, so that
\(d \alpha / h_1=l d s, \cdots \). Also let \(d s(1+e)\) be the length of
the corresponding linear element in the strained state. Then \(e\) is
given by the equation
\[
\begin{aligned}
{(1+e)}^2= & {\left[l\left\{1+h_1 \frac{\partial u_a}{\partial \alpha}+h_1 h_2
u_\beta \frac{\partial}{\partial \beta}\left(\frac{1}{h_1}\right)+h_1 h_3
u_\gamma \frac{\partial}{\partial
\gamma}\left(\frac{1}{h_1}\right)\right\}\right.} \\
& \left.+m \frac{h_2}{h_1} \frac{\partial}{\partial \beta}\left(h_1
u_a\right)+n \frac{h_3}{h_1} \frac{\partial}{\partial \gamma}\left(h_1
u_a\right)\right]^2+\cdots+\ldots .
\end{aligned}
\]
Neglecting squares and products of \(u_\alpha, u_\beta, u_\gamma \), we
may write the result in the form
\[
\varepsilon =
\begin{bmatrix}
e_{\alpha \alpha} & \frac{1}{2} e_{\alpha \beta} & \frac{1}{2} e_{\gamma \alpha}  \\
\frac{1}{2} e_{\alpha \beta}  & e_{\beta \beta} & \frac{1}{2} e_{\beta \gamma}  \\
\frac{1}{2} e_{\gamma \alpha}  & \frac{1}{2} e_{\beta \gamma} & e_{\gamma \gamma}
\end{bmatrix}, \quad
e=\begin{bmatrix} \ell & m & n \end{bmatrix}^T
\varepsilon \begin{bmatrix} \ell & m & n \end{bmatrix},
\]
in which
\[
\left.\begin{array}{l}
e_{a a}=h_1 \frac{\partial u_a}{\partial \alpha}+h_1 h_2 u_\beta
\frac{\partial}{\partial \beta}\left(\frac{1}{h_1}\right)+h_3 h_1 u_\gamma
\frac{\partial}{\partial \gamma}\left(\frac{1}{h_1}\right), \\
e_{\beta \beta}=h_2 \frac{\partial u_\beta}{\partial \beta}+h_2 h_3 u_\gamma
\frac{\partial}{\partial \gamma}\left(\frac{1}{h_2}\right)+h_1 h_2 u_\alpha
\frac{\partial}{\partial \alpha}\left(\frac{1}{h_2}\right), \\
e_{\gamma \gamma}=h_3 \frac{\partial u_\gamma}{\partial \gamma}+h_3 h_1
u_\alpha \frac{\partial}{\partial \alpha}\left(\frac{1}{h_3}\right)+h_2 h_3
u_\beta \frac{\partial}{\partial \beta}\left(\frac{1}{h_2}\right), \\
e_{\beta \gamma}=\frac{h_2}{h_3} \frac{\partial}{\partial \beta}\left(h_3
u_\gamma \right)+\frac{h_3}{h_2} \frac{\partial}{\partial \gamma}\left(h_2
u_\beta \right), \\
e_{\gamma \alpha}=\frac{h_3}{h_1} \frac{\partial}{\partial \gamma}\left(h_1
u_\alpha \right)+\frac{h_1}{h_3} \frac{\partial}{\partial \alpha}\left(h_3
u_\gamma \right), \\
e_{\alpha \beta}=\frac{h_1}{h_2} \frac{\partial}{\partial \alpha}\left(h_2
u_\beta \right)+\frac{h_2}{h_1} \frac{\partial}{\partial \beta}\left(h_1
u_\alpha \right) .
\end{array}\right\}
\]
\(\varepsilon \) is the strain referred to the orthogonal coordinates.
In fact
\(e_{\alpha \alpha}\) is the extension of a linear element which, in the
unstrained state, lies along the normal to the surface \(\alpha \); and
\(e_{\beta \gamma}\) is the cosine of the angle between the linear
elements which, in the unstrained state, lie along the normals to the
surfaces \(\beta \) and \(\gamma \).

\section[Curvilinear Rotation]{Dilatation and Rotation referred to curvilinear orthogonal coordinates.}
%\section{Dilatation and Rotation referred to curvilinear orthogonal coordinates.}% 21

The results of Art. 15 can be utilized to express the cubical dilatation
\(\Delta \), and the component rotations \(\varpi_a, \varpi_\beta, \varpi_\gamma \)
about the normals to the three surfaces, in terms of the components
\(u_\alpha, u_\beta, u_\gamma \) of the displacement.

To obtain the expression for \(\Delta \) we form the surface integral of
the normal component of the displacement\footnote{This method is due to Lord
Kelvin. (Sir W. Thomson, Math.\ and Phys.  Papers, Vol. 1, p. 25. The date of
the investigation is 1843.) }over the surface of an element of the body bounded
by the three pairs of surfaces
\((\alpha, \alpha+d \alpha),(\beta, \beta+d \beta),(\gamma, \gamma+d \gamma)\),
the normal being drawn away from the interior of the element. The
contributions of the faces of the element can be put down in such forms as
\[
\begin{aligned}
& \text { contribution of } \quad a=-u_a \frac{d \beta}{h_2} \frac{d \gamma}{h_3} \\
& \quad \text { " " } \quad a+d \alpha=u_a \frac{d \beta}{h_2}
\frac{d \gamma}{h_3}+d \alpha \frac{\partial}{\partial \alpha}\left(u_a \frac{d
\beta}{h_2} \frac{d \gamma}{h_3}\right),
\end{aligned}
\]
and, on adding the six contributions, we obtain
\[
d a d \beta d \gamma \left\{\frac{\partial}{\partial \alpha}
\left(\frac{u_a}{h_2 h_3}\right)+
\frac{\partial}{\partial \beta}\left(\frac{u_\beta}{h_3 h_1}\right)+
\frac{\partial}{\partial \gamma}\left(\frac{u_\gamma}{h_1 h_2}\right)\right\};
\]
this must be the same as \(\Delta d a d \beta d \gamma / h_1 h_2 h_2\).
We therefore have
\[
\Delta=h_1 h_2 h_3\left\{
\frac{\partial}{\partial \alpha}\left(\frac{u_a}{h_2 h_3}\right)+
\frac{\partial}{\partial \beta}\left(\frac{u_\beta}{h_3 h_1}\right)+
\frac{\partial}{\partial \gamma}\left(\frac{u_\gamma}{h_1 h_2}\right)
\right\}.
\]

This result is the same as would be found by adding the expressions for
\(e_{\alpha \alpha}\), \(e_{\beta \beta}\), \(e_{\gamma \gamma}\) in (36).
